\chapter{Diagrams as Data: S-Complexes}
\label{ch:qtsl-scomplexes}

\noindent\draftmark\medskip

Several strata of this book traffic in \emph{diagrams}. The
specification language's global stratum asserts a certified Morse graph
(\bookref{ch:qtsl-specification-language}{the Specification Language chapter}); the hypernet
chapter reads every computation as a string diagram
(\bookref{ch:qtsl-hypernets}{the hypernet chapter}); grid
certificates decompose state space into cubical cells; and the
higher-categorical diagram machinery this project consults runs on
oriented graded posets. Each of these is, concretely, a family of
cells in graded dimensions plus incidence data saying which cells
bound which. Each could therefore grow its own bespoke incidence
schema, its own well-formedness checker, and its own subtly different
answer to ``is the boundary recorded here the boundary the theorem
was about?''

The toolkit owns one such object. An \emph{S-complex} is a graded
family of finite cell sets whose boundary structure is recorded as a
pair of semiring-valued tensors per dimension (one for each
orientation sign), satisfying a subtraction-free form of the chain
condition $\partial^2 = 0$. The semiring parameter $\mathbb{S}$ is where
the genus gets its reach: with Booleans it is incidence, with natural
numbers it is multiplicity, with a semiring of \emph{port names} it
is a string diagram, and with min-plus dwell bounds its labels carry
quantitative dynamics. Well-formedness in every case amounts to a
handful of tensor contractions and one semiring comparison per
dimension, using the same contraction machinery the rest of the
toolkit already trusts (\bookref{ch:qtsl-hypernets}{the hypernet
chapter}, relational-kinship section), so the properties are
checkable in bulk, on accelerators, by the engine that checks
everything else.

The instinct (diagrams stored as typed incidence data, computed on
in bulk) is shared with the attributed $\mathcal{C}$-set line of
applied category theory, where diagrams are categorical data
structures~\cite{patterson21} with data-parallel composition
algorithms~\cite{wilson23}, and where a string diagram can be a
first-class engineering artefact rather than an
illustration~\cite{master19}. An S-complex trades that line's schema
generality for one graded shape with a semiring parameter; the trade
is what makes well-formedness a tensor computation.

This chapter defines the object and its morphisms
(section~\ref{sec:scx-def}), tours the instances that choice of
$\mathbb{S}$ buys (section~\ref{sec:scx-instances}), states the
encoding discipline that decides when a complex is stored versus
derived (section~\ref{sec:scx-generators}), presents the certified
Morse graph as the flagship rank-1 instance with its quantitative
refinement (section~\ref{sec:scx-morse}), works the string-diagram
instance in full --- including the reading under which a well-formed
hypernet \emph{is} an S-complex (section~\ref{sec:scx-strings}) ---
and closes with a ledger separating definitional content from design
directions (section~\ref{sec:scx-ledger}).
Sections~\ref{sec:scx-def} and~\ref{sec:scx-morse} carry everything
the rest of the book consumes; section~\ref{sec:scx-strings} is the
most likely to surprise, and a reader with no stake in string
diagrams can skip it without loss of continuity.

\section{The definition}
\label{sec:scx-def}

Throughout, $\mathbb{S}$ is an \emph{ordered semiring}: a semiring
$(\mathbb{S}, 0, 1, +, \cdot)$ equipped with a preorder $\le$ such
that $0 \le 1$ and both operations are monotone in each argument.
Order-equivalence is derived, not assumed:
$x \cong y \mathrel{:\Leftrightarrow} x \le y \wedge y \le x$.
No subtraction is ever required. That restraint is the point: the
instances worth having --- Booleans, counts, port sets, min-plus
costs --- live in semirings that have no subtraction to offer.

A \emph{tensor} over $\mathbb{S}$ indexed by finite sets
$A_0, \dots, A_{k-1}$ is a finite function
$\prod_{i<k} A_i \to \mathbb{S}$, and juxtaposition of tensors
denotes contraction over the shared index (sum of products, in
$\mathbb{S}$), as in the contraction vocabulary of
\bookref{ch:qtsl-mathematical-models}{the modelling chapter}.

\begin{definition}[S-complex]
\label{def:scomplex}
An $\mathbb{S}$-\emph{complex} of rank $r$ consists of finite cell
sets $C_0, C_1, \dots, C_r$ (cells in dimension $0$ through $r$)
together with, for each $0 < i \le r$ and each sign
$e \in \{+,-\}$, a rank-2 boundary tensor
\[
\partial_i^{e} \colon C_{i} \times C_{i-1} \longrightarrow \mathbb{S},
\]
satisfying, for each $1 \le i \le r-1$, the \emph{sign-split chain
condition}
\begin{equation}
\label{eq:scx-dd}
\partial_{i}^{+}\, \partial_{i+1}^{+} \;+\; \partial_{i}^{-}\, \partial_{i+1}^{-}
\;\cong\;
\partial_{i}^{-}\, \partial_{i+1}^{+} \;+\; \partial_{i}^{+}\, \partial_{i+1}^{-}
\qquad \colon\; C_{i+1} \times C_{i-1} \to \mathbb{S}.
\end{equation}
\end{definition}

Stated as one dependent type, with $[N]$ the finite ordinals and
$\mathsf{Fin}$ the concretely finite sets, the presented data of the
genus is
\[
\sum_{N \colon \mathbb{N}} \;\; \sum_{C \colon \prod_{n \colon [N]} \mathsf{Fin}} \;\; \prod_{i \colon [N-1]} \; \prod_{e \colon \{+,-\}} \big[\, \partial_{i+1}^{e} \colon C_{i+1} \times C_{i} \to \mathbb{S} \,\big]
\]
--- $N$ graded levels and a sign-split boundary pair per consecutive
pair of levels, nothing else. The type records presentations only;
condition~\eqref{eq:scx-dd} is a property checked \emph{of} a
presentation rather than a component packaged \emph{in} it, so one
schema serves both the complexes that satisfy it and the raw
presentations a checker has not yet admitted. The rank $r$ of
Definition~\ref{def:scomplex} is $N - 1$.

Condition~\eqref{eq:scx-dd} is the classical $\partial^2 = 0$ with
the subtraction moved across the equals sign: over a ring, setting
$\partial = \partial^{+} - \partial^{-}$ and expanding
$\partial_i\,\partial_{i+1} = 0$ gives precisely ``matched sign pairs
equal mismatched sign pairs.'' Stated this way the condition makes
sense over any ordered semiring, and each side is two contractions
and a sum, so checking well-formedness of a presented complex is a
fixed, small pattern of contraction rows per dimension --- the
bulk-checkable shape the engine is built around. An elementary
instance to hold while reading: a filled square. Each corner lies on
two of the four edges, so exactly two routes descend from the
$2$-cell through an edge to that corner; however the edges are
oriented, one route carries matched signs and the other mismatched,
so the two sides of~\eqref{eq:scx-dd} agree entry by entry --- the
corner-cancellation of $\partial^2 = 0$, surviving with parity where
its minus signs used to be. For $r \le 1$ the
condition is vacuous, and a presentation of a rank-1 complex should
say so plainly rather than imply a check that never ran.

The two signs record \emph{combinatorial orientation}: for a directed
edge $a$ from $x$ to $y$, the source incidence lives in
$\partial_1^{-}(a, x)$ and the target incidence in
$\partial_1^{+}(a, y)$. This sign records a property of the diagram
itself rather than of any dynamics drawn on it; when a dynamical
certificate needs a notion of inward and outward flow across a face,
that is separate label data riding on the cells, never a
reinterpretation of the boundary signs.

\begin{definition}[morphism of S-complexes]
\label{def:scx-morphism}
Let $C$ be an $\mathbb{S}$-complex and $D$ a $\mathbb{T}$-complex
with $\operatorname{rank}(C) \le \operatorname{rank}(D)$. A
\emph{morphism} $f \colon C \to D$ consists of a function
$f_i \colon C_i \to D_i$ for each dimension $i$ of $C$, regarded
semantically as a tensor
$f_i \colon C_i \times D_i \to \mathbb{T}$ whose entries are $0$
and $1$, together with a lax
morphism of ordered semirings
$f_S \colon \mathbb{S} \to \mathbb{T}$ (monotone, $f_S(0) \cong 0$,
$f_S(1) \cong 1$, $f_S(x+y) \le f_S(x) + f_S(y)$, and
$f_S(x \cdot y) \le f_S(x) \cdot f_S(y)$), such that for every
dimension $0 < i \le \operatorname{rank}(C)$, every sign $e$, and
every $(c, d) \in C_i \times D_{i-1}$,
\begin{equation}
\label{eq:scx-chainmap}
\sum_{d' \in D_i} f_i(c, d')\; \partial_i^{D,e}(d', d)
\;=\;
\sum_{c' \in C_{i-1}} f_S\!\big(\partial_i^{C,e}(c, c')\big)\; f_{i-1}(c', d).
\end{equation}
\end{definition}

Equation~\eqref{eq:scx-chainmap} is the chain-map square, sign by
sign: mapping a cell forward and taking its boundary in $D$ agrees
with taking its boundary in $C$, pushing the cells forward, and
translating the coefficient values through $f_S$. Laxity in $f_S$ is
deliberate. The useful coefficient translations are
\emph{forgetful} --- cardinality (from port witnesses to counts),
support (from counts to incidence), and coarse quantitative summaries
(from fine margins to dwell bounds) --- and translations of the last
kind are subadditive rather than exact, so demanding exactness would
refuse precisely the summaries the tower is for.

\section{Instances by choice of semiring}
\label{sec:scx-instances}

The definition specializes by choice of coefficient semiring:

\begin{center}
\begin{tabular}{@{}llll@{}}
\toprule
object & $\mathbb{S}$ & rank & boundary reading \\
\midrule
directed graph & $\mathbf{2}$ (Booleans) & 1 & $\partial_1^{-}$ source, $\partial_1^{+}$ target \\
directed hypergraph & $\mathbb{N}$ & 1 & tail/head multiplicities \\
cubical complex & $\mathbf{2}$ & $d$ & oriented facet incidence \\
oriented graded poset & $\mathbf{2}$ & any & covering relation, one sign each \\
Lefschetz complex & a ring's split halves & any & incidence coefficients \\
string diagram & port names (\S\ref{sec:scx-strings}) & 1 & \emph{which ports}, not just how many \\
\bottomrule
\end{tabular}
\end{center}

\begin{remark}[oriented graded posets]
\label{rem:scx-ogp}
Oriented graded posets, a combinatorial substrate of
higher-categorical diagram rewriting \cite{hadzihasanovic24a}, embed
as Boolean complexes in which each covering pair carries exactly one
of the two signs. The shape predicates that theory cares about
(roundness, molecule-hood, regularity) remain predicates on top of
the complex, as that literature states them, with the complex
supplying the data they are predicates \emph{of}.
\end{remark}

\begin{remark}[Lefschetz complexes]
\label{rem:scx-lefschetz}
Lefschetz complexes, a combinatorial carrier of Conley-theoretic
invariants \cite{mischaikowmrozek2002}, carry ring coefficients
without a ring: a classical incidence coefficient enters as its
split halves $k = k^{+} - k^{-}$, one sign each, so entries stay in
$\mathbb{N}$ and subtraction never enters the schema.
\end{remark}

\begin{remark}[the coefficient tower]
\label{rem:scx-tower}
The coefficient levels form a tower. Cardinality maps a string
diagram's witness sets to a directed hypergraph's multiplicities, and
Boolean support forgets those to bare incidence. Both maps are lax
morphisms in the sense of Definition~\ref{def:scx-morphism} with the
identity on cells, so the port, counting, and incidence readings are
three resolutions of a single object.
\end{remark}

\section{Generators or presentations}
\label{sec:scx-generators}

Not every complex the toolkit reasons about should be stored as one.
The encoding discipline distinguishes two postures.

\begin{itemize}
\item \textbf{Generator-encoded} where the complex is canonically
  derivable: grid certificates store grid lines and cell selections,
  never cell lists or incidence tensors --- the cubical complex, with
  its orientation-split boundary, is recomputed from the generators
  by every consumer. Skew between a stored presentation and its
  generators is then unrepresentable rather than checked.
\item \textbf{S-complex-encoded} where the boundary \emph{is} the
  content: a Morse graph's edges, a reduced complex, a connection
  matrix. Here there is no generator to fall back on --- the
  incidence data is precisely what the certificate asserts --- so the
  presentation of Definition~\ref{def:scomplex} is the wire object,
  and condition~\eqref{eq:scx-dd} (where non-vacuous) is part of its
  well-formedness check.
\end{itemize}

The two postures compose: a quotient from a generator-encoded complex
onto an S-complex-encoded one is a morphism in the sense of
Definition~\ref{def:scx-morphism}, with the derived presentation
appearing only transiently on the checking side.

\section{The Morse graph, quantitatively}
\label{sec:scx-morse}

The global stratum's payload
(\bookref{ch:qtsl-specification-language}{the Specification Language chapter}, QMorse row) is
a certified Morse graph: finitely many blocks, finitely many directed
connections. As data it ships in adjacency dialect (a vertex count
and an edge list), and that dialect \emph{is} the rank-1 Boolean
S-complex: $C_0$ the Morse sets, $C_1$ the connections,
$\partial_1^{-}$ source, $\partial_1^{+}$ target. Acyclicity is the
checker's to confirm; the confirmation, in this vocabulary, is an
order witness on $C_0$ of the same kind
section~\ref{sec:scx-strings} uses for string diagrams.

The quantitative refinement adds a \emph{labelling} without changing
the complex: cells carry values in the ordered semiring
$\mathrm{DWELL} = ([0, \infty], \min, +)$ --- the cost quantale of
Lawvere's generalized metric spaces~\cite{lawvere73}, whose additive
unit is $\infty$ and whose multiplicative unit is $0$, ordered by
the reverse of the numerical order, which is what makes $0 \le 1$
hold --- with
connections labelled by certified lower bounds on transit time and
Morse sets labelled by their sojourn bounds, which for genuine Morse
sets are $\infty$. The semiring arithmetic then guards the composite
claims against overstatement: min-plus contraction of a path through
an intermediate Morse set yields $\infty$, so the algebra itself
reports that no finite total-time bound through a recurrent set is
being asserted, while per-leg bounds stay finite and usable. This is
the intended serialization form for the quantitative Morse tier: the
graph presented as in Definition~\ref{def:scomplex}, the rates as
cell labels (extended nonnegative rationals in the machine
artefact), and the projection from the underlying grid complex
onto the Morse graph stated as a morphism --- cells to their Morse
or transit class, fine certified margins translated laxly to dwell
bounds --- rather than as an unstated convention between two
schemas.

\section{String diagrams: the port semiring}
\label{sec:scx-strings}

This instance shows why the coefficients form a \emph{semiring
parameter} rather than a Boolean convention. Fix a supply
$\mathcal{P}$ of \emph{port names} --- opaque identifiers, meaningful
only up to renaming. Let $\mathbb{S}_{\mathrm{port}}$ be the semiring
of finite \emph{multisets} of finite sequences over $\mathcal{P}$:
addition is multiset union with zero $\varnothing$, multiplication is
elementwise concatenation with unit $\{\varepsilon\}$, and the order
compares cardinalities --- so the derived order-equivalence is blind
to the names themselves, which is the renaming-invariance the names'
opacity demands.\footnote{Multisets rather than sets, for a reason
worth checking: concatenation can merge distinct pairs of sequences
($\langle 1\rangle \cdot \langle 2,3\rangle$ and
$\langle 1,2\rangle \cdot \langle 3\rangle$ both concatenate to
$\langle 1,2,3\rangle$), so two \emph{sets} of equal cardinality can
have products of different cardinalities, and multiplication would
fail the monotonicity section~\ref{sec:scx-def} demands of every
coefficient semiring. Multisets count collisions instead of merging
them, which makes cardinality exactly additive and multiplicative.
The witness sets arising in diagrams are unaffected: distinct ports
give distinct sequences.}

A \emph{string diagram} over a signature is a rank-1
$\mathbb{S}_{\mathrm{port}}$-complex --- $C_0$ the wires, $C_1$ the
boxes --- together with a morphism into the complex presenting the
signature ($C_0$ the sorts, $C_1$ the generators, each generator's
canonical ports as its boundary values). A boundary value
$\partial_1^{-}(a, x)$ is a \emph{witness set}: one singleton
sequence per input port of box $a$ attached to wire $x$, and dually
for outputs. Over $\mathbf{2}$ an edge could only assert incidence;
over $\mathbb{S}_{\mathrm{port}}$ the value says \emph{what flows
where} --- the difference between a reachability structure and a
diagram whose wires carry content comes down, in this vocabulary, to
the choice of semiring. The signature morphism does the typing: its
cell maps send boxes to generators and wires to sorts, and its
coefficient translation $f_S$ matches each diagram port to the
generator port it instantiates, so argument order lives in the
signature once instead of in every diagram. Sequence concatenation
earns its place under composition, where contracted products
accumulate port \emph{paths} witnessing composite incidences; at
rank 1 only singletons occur. Figure~\ref{fig:scx-tower} works a
two-box example and reads its boundary data at the three coefficient
resolutions of Remark~\ref{rem:scx-tower}.

\begin{figure}[t]
\centering
\begin{tikzpicture}[x=0.94cm,y=0.82cm,baseline=(current bounding box.center),
  box/.style={draw=accentdeep, fill=accent!50, rounded corners=1pt,
    font=\footnotesize\ttfamily, inner sep=2.5pt, minimum height=4.2mm},
  wire/.style={-, line width=0.5pt, draw=black!75},
  port/.style={font=\small\itshape, inner sep=1pt},
  plab/.style={font=\scriptsize, inner sep=1pt, text=black!60}]
\node[port] (x) at (0, 0.75) {$x$};
\node[box]  (a) at (1.4, 0.75) {$a$};
\node[box]  (b) at (4.4, 0.75) {$b$};
\node[port] (z) at (5.8, 0.75) {$z$};
\coordinate (split) at (2.5, 0.75);
\draw[wire] (x) -- (a);
\draw[wire] (a.east) -- (split);
\draw[wire] (split) .. controls (3.1,0.75) and (3.5,1.15) ..
  ([yshift=1.6mm]b.west);
\draw[wire] (split) .. controls (3.1,0.75) and (3.5,0.35) ..
  ([yshift=-1.6mm]b.west);
\draw[wire] (b) -- (z);
\node[plab, anchor=south] at (2.15, 0.85) {$y$};
\node[plab, anchor=south] at ([xshift=-2.6mm,yshift=1.5mm]b.west) {1};
\node[plab, anchor=north] at ([xshift=-2.6mm,yshift=-1.5mm]b.west) {2};
\end{tikzpicture}
\hspace{1.2em}
\begin{tabular}{@{}lccc@{}}
\toprule
entry & $\mathbb{S}_{\mathrm{port}}$ & $\mathbb{N}$ & $\mathbf{2}$ \\
\midrule
$\partial_1^{-}(a,x)$ & $\{\langle 1\rangle\}$ & $1$ & $1$ \\
$\partial_1^{+}(a,y)$ & $\{\langle 1\rangle\}$ & $1$ & $1$ \\
$\partial_1^{-}(b,y)$ & $\{\langle 1\rangle,\langle 2\rangle\}$ & $2$ & $1$ \\
$\partial_1^{+}(b,z)$ & $\{\langle 1\rangle\}$ & $1$ & $1$ \\
\bottomrule
\end{tabular}
\caption{One diagram, three coefficient resolutions. Wire $y$ feeds
both input ports of box $b$, so the witness set
$\partial_1^{-}(b,y)$ holds two singleton port sequences; cardinality
forgets \emph{which ports} to \emph{how many}, and Boolean support
forgets that to \emph{whether at all} (Remark~\ref{rem:scx-tower}).
Entries not shown are $\varnothing$, $0$, $0$.}
\label{fig:scx-tower}
\end{figure}

Structural conditions on the diagram are then one clause each, and
free-structure theorems select which clauses to impose:

\begin{center}
\begin{tabular}{@{}lll@{}}
\toprule
producer/consumer discipline & acyclicity & free structure presented \\
\midrule
none & none & hypergraph category \cite{bonchi20} \\
$\le 1$ producer, $\le 1$ consumer & none & traced symmetric monoidal \cite{kissinger13} \\
$\le 1$ producer, $\le 1$ consumer & witness & symmetric monoidal \cite{bonchi22} \\
$\le 1$ producer, consumers free & witness & gs-monoidal \cite{fritzliang22} \\
\bottomrule
\end{tabular}
\end{center}

Here ``$\le 1$ producer'' is the port-counting condition
$\sum_{a \in C_1} \lvert \partial_1^{+}(a, x) \rvert \le 1$ for every
wire $x$, and dually for consumers --- the cardinality order makes
counting a first-class operation. The acyclicity witness is a
function $t \colon C_0 \to P$ into a poset with
$t(\text{input wire}) < t(\text{output wire})$ across every box: a
one-pass check and a rewrite-friendly certificate.

\subsection{Hypernets read as S-complexes}
\label{sec:scx-hypernets}

The hypernet of \bookref{ch:qtsl-hypernets}{the hypernet chapter}
sits in the third row of the table: its linearity condition is
producer-and-consumer monogamy and its static-single-assignment
order is an acyclicity witness, as that chapter already observes in
citing the free-structure correspondence. The reading assembles as
follows. Wires are the 0-cells; applications are the 1-cells; an
application's ordered input and output lists are the two boundary
rows, with list position serving as the port name --- the
representation takes the renaming quotient so the semiring does not
have to. The catalogue plays the signature: the morphism sends each
application to its box, each wire to its tensor type, and its
coefficient translation records which argument slot each wire fills.
Two strata of the hypernet deliberately live \emph{beside} the
complex rather than in it: box payloads (constants, contraction
patterns, polynomial coefficients) are cell labels, just as rates
are labels on the Morse graph; and the interface is a cospan of
morphisms from rank-0 complexes --- an ordered list of graph inputs
and one of graph outputs, each simply a function into $C_0$, which
is what the declared interface lists already are.

The vocabulary also makes the design space around the interface
legible in single clauses. Letting the output leg of the cospan hit
any wire --- including one an application already consumes ---
declares a read-out port without touching the inner complex or its
monogamy, since the legs were never part of either. If consumer-side
monogamy is retired while the producer side is kept, the diagram
genus moves one row down the table, to the gs-monoidal discipline
of every static-single-assignment compiler representation --- def
once, use many --- with copy and discard absorbed into the wiring.
Both are stated here as readings the definition supports rather than
as changes made: the shipped hypernet remains the monogamous, acyclic
row, and its bytes, digests, and substitution laws are those of
\bookref{ch:qtsl-hypernets}{the hypernet chapter}, unchanged.

\section{Ledger}
\label{sec:scx-ledger}

What in this chapter is what:

\begin{itemize}
\item \textbf{Definitional.} Definitions~\ref{def:scomplex}
  and~\ref{def:scx-morphism}, the derivation of~\eqref{eq:scx-dd}
  from $\partial^2 = 0$, the instance table, and the free-structure
  table (whose correspondences are the cited literature's theorems,
  not this book's).
\item \textbf{Machine artefact.} The genus is declared as a
  standard-library theory of the geolog engine
  (\path{theories/SComplex.geolog}), which names this chapter as its
  expository companion, and the chain condition's contraction
  factoring is generated from a surface statement of the tensor law
  (\path{theories/SComplexGenus.geolog}). The theory is staged
  design vocabulary, exercised by the engine's tests rather than by
  production data.
\item \textbf{Shipped elsewhere in this book.} The Morse-graph
  payload shape and its checker obligations
  (\bookref{ch:qtsl-specification-language}{the Specification Language chapter}); the
  hypernet IR, whose linearity and ordering conditions this
  chapter's section~\ref{sec:scx-hypernets} re-reads
  (\bookref{ch:qtsl-hypernets}{the hypernet chapter}); the
  contraction machinery that makes
  condition~\eqref{eq:scx-dd} bulk-checkable
  (\bookref{ch:qtsl-mathematical-models}{the modelling chapter}).
\item \textbf{Design tense.} The quantitative Morse serialization in
  complex-plus-labels form and its quotient-as-morphism statement
  (section~\ref{sec:scx-morse}); any S-complex-native storage for
  hypernets (section~\ref{sec:scx-hypernets} is a reading, and the
  canonical bytes are unchanged by it); mechanized S-complex
  morphisms, of which the toolkit currently has none --- the Morse
  quotient is the natural first target. Higher-dimensional
  rewriting on this substrate (equivalence cells over diagrams) is
  out of this chapter's scope entirely.
\end{itemize}
